\documentclass[11pt]{article}

\usepackage[margin=1in]{geometry}
\usepackage[T1]{fontenc}
\usepackage[utf8]{inputenc}
\usepackage{lmodern}
\usepackage{hyperref}
\usepackage{enumitem}

\title{Do Self-Supervised Vision Models Learn What Experts See?\\
\large Report Skeleton for the Final Project Write-Up}
\author{Team Name / Member Names}
\date{\today}

\begin{document}

\maketitle

% NOTE:
% This is a planning skeleton intended to be easy to paste into the official
% course Overleaf template later.
%
% If you are using the course template:
% 1. Copy the body content from \begin{abstract} onward into the template.
% 2. Replace the title/author block with the template's required fields.
% 3. Replace bullet placeholders with final prose, figures, and tables.

\begin{abstract}
This abstract should briefly state the problem, the dataset, the study design,
and the headline findings.

Suggested content:
\begin{itemize}[leftmargin=1.5em]
    \item One sentence on the core problem: whether SSL vision models attend to
    the same diagnostic regions that experts mark as important.
    \item One sentence on the setting: WikiChurches, expert bounding boxes, and
    architectural feature analysis.
    \item One sentence on the three-question study design: frozen-model
    alignment, fine-tuning shift, and per-head specialization.
    \item One sentence summarizing the main findings once finalized.
\end{itemize}
\end{abstract}

\section{Project-at-a-Glance Overview}

This section can be a short opening infographic or a compact summary block.

\begin{itemize}[leftmargin=1.5em]
    \item Number of models compared.
    \item Number of fine-tuning strategies.
    \item Number of alignment metrics.
    \item Number of annotated images and expert bounding boxes.
    \item Number of research questions.
    \item One short sentence explaining what the study contributes overall.
\end{itemize}

\section{Introduction and Motivation}

\begin{itemize}[leftmargin=1.5em]
    \item Explain why classification accuracy alone is not enough.
    \item Motivate the difference between being correct and attending to the
    right visual evidence.
    \item Explain why expert-annotated architectural features make this a
    useful testbed.
    \item End with the main gap addressed by the project:
    \begin{itemize}[leftmargin=1.5em]
        \item cross-model expert-alignment comparison,
        \item post-fine-tuning attention-shift analysis,
        \item per-head specialization analysis.
    \end{itemize}
\end{itemize}

\section{Research Questions and Contributions}

\subsection{Q1: Frozen-Model Attention Alignment}

\begin{itemize}[leftmargin=1.5em]
    \item State the question clearly: do frozen SSL and baseline models attend
    to expert-marked diagnostic features?
    \item Frame this as the core benchmark question.
\end{itemize}

\subsection{Q2: Fine-Tuning and Attention Shift}

\begin{itemize}[leftmargin=1.5em]
    \item Ask how attention changes after adaptation.
    \item Compare Linear Probe, LoRA, and Full fine-tuning.
    \item Clarify that this question is about both direction and magnitude of
    shift.
\end{itemize}

\subsection{Q3: Per-Head Specialization}

\begin{itemize}[leftmargin=1.5em]
    \item Ask whether individual heads specialize for different feature types.
    \item Keep the scoped study claim clear and modest.
\end{itemize}

\subsection{Contributions}

\begin{itemize}[leftmargin=1.5em]
    \item Multi-metric benchmark for expert-aligned attention evaluation.
    \item Calibrated Q1 comparison against naive baselines.
    \item Q2 fine-tuning analysis across three adaptation strategies.
    \item Q3 scoped per-head specialization study.
    \item Reproducible pipeline and interactive analysis tool.
\end{itemize}

\section{Related Work}

\begin{itemize}[leftmargin=1.5em]
    \item Attention interpretability for vision transformers.
    \item Evaluation against human or expert annotations.
    \item Fine-tuning, attention drift, and representational change.
    \item Attention-head specialization.
    \item Cultural heritage / architectural recognition context.
    \item Position this project as a domain-grounded evaluation study rather
    than a new interpretability algorithm.
\end{itemize}

\section{Dataset and Problem Setup}

\begin{itemize}[leftmargin=1.5em]
    \item Introduce WikiChurches and justify its use.
    \item Describe the annotated subset and the style-labeled training pool.
    \item Explain the role of the expert bounding boxes.
    \item Summarize preprocessing or filtering decisions.
    \item State known scope limits such as sparse annotation bias.
\end{itemize}

\section{Methodology}

\subsection{Models and Attention Extraction}

\begin{itemize}[leftmargin=1.5em]
    \item Briefly describe the compared model families.
    \item Explain which attention-extraction method is used for each model.
    \item Distinguish CLS attention, rollout, mean attention, and Grad-CAM.
\end{itemize}

\subsection{Alignment Metrics}

\begin{itemize}[leftmargin=1.5em]
    \item Explain the role of each metric at a high level.
    \item Distinguish threshold-based metrics from threshold-free metrics.
    \item Make metric direction explicit:
    \begin{itemize}[leftmargin=1.5em]
        \item higher is better for IoU and Coverage,
        \item lower is better for MSE, KL, and EMD.
    \end{itemize}
\end{itemize}

\subsection{Baselines and Calibration}

\begin{itemize}[leftmargin=1.5em]
    \item Explain why continuous metrics need calibration.
    \item Introduce the naive baselines used for Q1 interpretation.
    \item State what these baselines help the reader understand.
\end{itemize}

\subsection{Fine-Tuning Protocol}

\begin{itemize}[leftmargin=1.5em]
    \item Describe the training, validation, and evaluation split logic.
    \item State clearly that the annotated evaluation images are excluded from
    the primary train/validation split.
    \item Summarize the three strategies:
    \begin{itemize}[leftmargin=1.5em]
        \item Linear Probe,
        \item LoRA,
        \item Full fine-tuning.
    \end{itemize}
    \item Explain checkpoint selection at a high level.
\end{itemize}

\subsection{Q3 Per-Head Scope}

\begin{itemize}[leftmargin=1.5em]
    \item Define the narrower scoped study for per-head analysis.
    \item State which models and variants support the primary Q3 claim.
    \item Keep methodological and interpretive guardrails explicit.
\end{itemize}

\subsection{Statistical Analysis}

\begin{itemize}[leftmargin=1.5em]
    \item Briefly describe the paired testing approach.
    \item Mention multiple-comparison correction.
    \item Mention effect sizes and confidence intervals.
\end{itemize}

\subsection{Methodological Safeguards and Reproducibility}

\begin{itemize}[leftmargin=1.5em]
    \item Summarize the choices that improve trustworthiness of the results.
    \item Suggested points:
    \begin{itemize}[leftmargin=1.5em]
        \item deterministic or stable analysis paths where needed,
        \item holdout discipline,
        \item calibrated baselines,
        \item explicit experiment provenance,
        \item artifact-based reporting workflow.
    \end{itemize}
\end{itemize}

\section{System and Analysis Interface}

\begin{itemize}[leftmargin=1.5em]
    \item Keep this section shorter than Methodology and Results.
    \item Explain that the software system enables the study rather than being
    the sole contribution.
    \item Summarize only the major components:
    \begin{itemize}[leftmargin=1.5em]
        \item metric and cache pipeline,
        \item fine-tuning experiment pipeline,
        \item backend and frontend analysis surfaces.
    \end{itemize}
    \item Avoid route-by-route API detail in the main report.
\end{itemize}

\section{Results}

\subsection{Q1 Results: Frozen-Model Attention Alignment}

\begin{itemize}[leftmargin=1.5em]
    \item Present the main frozen-model comparison.
    \item Include calibrated interpretation against baselines.
    \item Discuss cross-metric agreement and disagreement.
    \item Use a small number of representative examples.
\end{itemize}

\subsection{Q2 Results: Fine-Tuning Effects on Attention}

\begin{itemize}[leftmargin=1.5em]
    \item Make this one of the larger result subsections.
    \item Compare shifts across strategies and model families.
    \item Discuss which models improve, preserve, or degrade in alignment.
    \item Include one or two image-level examples such as shift-map
    interpretation if useful.
\end{itemize}

\subsection{Q3 Results: Per-Head Specialization}

\begin{itemize}[leftmargin=1.5em]
    \item Present head rankings, head-feature patterns, and exemplar
    observations.
    \item Keep claims descriptive unless stronger causal evidence is available.
    \item Keep this subsection scoped if the findings are less mature than Q1
    and Q2.
\end{itemize}

\section{Discussion}

\subsection{Intuitive vs Surprising Findings}

\begin{itemize}[leftmargin=1.5em]
    \item Highlight which outcomes matched expectations.
    \item Highlight which outcomes were surprising.
    \item Explain why those results may have occurred.
\end{itemize}

\subsection{What the Results Suggest About Pretraining Objectives}

\begin{itemize}[leftmargin=1.5em]
    \item Compare what different training paradigms appear to encourage.
    \item Relate observed behavior back to model families where appropriate.
\end{itemize}

\subsection{Practical Implications}

\begin{itemize}[leftmargin=1.5em]
    \item Explain what the findings could mean for domain adaptation and model
    selection.
    \item Keep this grounded in the actual study.
\end{itemize}

\subsection{Limitations and Threats to Validity}

\begin{itemize}[leftmargin=1.5em]
    \item Dataset size and domain scope.
    \item Sparse annotation bias.
    \item Limits of treating attention as explanation.
    \item Any limits in Q3 interpretation.
\end{itemize}

\section{Conclusion}

\begin{itemize}[leftmargin=1.5em]
    \item Restate the problem and the three research questions.
    \item Summarize the overall contribution of the study.
    \item End with a short future-work note only if needed.
\end{itemize}

\appendix

\section{Appendix}

\begin{itemize}[leftmargin=1.5em]
    \item Extra qualitative examples.
    \item Supplementary tables.
    \item Additional implementation detail.
    \item Experiment artifact references.
    \item Supplementary figures.
\end{itemize}

% Optional bibliography block once references are ready:
% \bibliographystyle{plain}
% \bibliography{references}

\end{document}
